\section{Independent orders, independent shifts, and every bilinear contact}
\label{ct:section}

The incoming \emph{Bilinear Tails and Determinants} report establishes
bilinear identities for two independently shifted alternating Hurwitz tails\cite{ProveItBilinear},
but leaves their higher contact hyperplanes explicitly open. Its diagonal
calculation does not determine mixed derivatives transverse to the diagonal.
We resolve that question by a finite polynomial calculus. In addition to
all cutoff finite parts, the calculation produces ordinary convergent
identities on every negative-total-order hyperplane. The first such family
has the particularly simple correction $1/2$ in every summand.

The finite bilinear identities recalled below are inherited. The new
statements are the centered coefficient algebra, the independent contact
and translation formulas, their sharp spectral degree, and the ordinary
identities extracted from them. The classical analytic inputs are the
Gamma quotient expansion and the alternating Hurwitz Euler expansion;
no assertion of global priority for these ingredients is intended.

\subsection{The inherited endpoint law and its normalization}

On $\Re z>0$ put
\begin{align}
 \mathcal B_s(z)&=2^{-s}\left\{\zeta\left(s,\frac z2\right)
              -\zeta\left(s,\frac{z+1}{2}\right)\right\},
                   \label{ct:B}\\
 h(z)&=\sqrt2\,\frac{\Gamma((z+1)/2)}{\Gamma(z/2)},
 \qquad \ell(z)=\log h(z),\qquad
 \tau_s(z)=e^{s\ell(z)}\mathcal B_s(z),\label{ct:tau}\\
 V_s(z)&=\frac{z^s-(z-1)^s}{h(z)^{2s}},\qquad \Re z>1.
                   \label{ct:V}
\end{align}
All powers use the holomorphic logarithms real on the positive axis.
The common Hurwitz pole in \eqref{ct:B} cancels, so this expression is
entire in $s$. The classical recurrences give
\begin{equation}\label{ct:recurrences}
 h(z)h(z+1)=z,\qquad
 \mathcal B_s(z)+\mathcal B_s(z+1)=z^{-s},\qquad
 \tau_s(z-1)-\tau_s(z+1)=V_s(z)\tau_s(z).
\end{equation}
The last identity follows by multiplying the two adjacent Hurwitz
recurrences by the appropriate Gamma powers and subtracting.

Let $\Re a,\Re b>0$ and let $p,q$ be independent complex orders. Define,
for $n\ge-1$,
\begin{align}
 u_n&=\tau_p(n+a+1),&v_n&=\tau_q(n+b+1),\notag\\
 C_n&=u_nv_{n+1}+u_{n+1}v_n,&
 D_n&=u_nv_{n+1}-u_{n+1}v_n,\label{ct:CD}\\
 E_+&=C_{-1},&E_-&=D_{-1}.\notag
\end{align}
For later convenience, the signs are kept outside the definitions
\[
 g_n^+=\{V_p(n+a+1)+V_q(n+b+1)\}u_nv_n,
 \qquad
 g_n^-=\{V_p(n+a+1)-V_q(n+b+1)\}u_nv_n.
\]
The source's exact finite identities are
\begin{equation}\label{ct:finite}
 S_N^+:=\sum_{n=0}^N g_n^+=E_+-C_N,
 \qquad
 S_N^-:=\sum_{n=0}^N(-1)^ng_n^-=E_-+(-1)^ND_N.
\end{equation}
Indeed, \eqref{ct:recurrences} gives
$g_n^+=C_{n-1}-C_n$ and $g_n^-=D_{n-1}+D_n$.
These identities hold at every complex order before taking a limit.
For $\Re(p+q)>0$ both infinite sums converge absolutely to $E_+$ and
$E_-$ respectively. The entire functions $E_\pm$ are their endpoint
continuations. A cutoff finite part need not equal that continuation.

\subsection{A centered finite coefficient algebra}

Use Bernoulli polynomials with $B_1(x)=x-1/2$ and Euler numbers normalized by
\[
 \operatorname{sech}t=\sum_{j\ge0}E_{2j}\frac{t^{2j}}{(2j)!}.
\]
Thus $E_0=1$, $E_2=-1$, and $E_4=5$. Put
\begin{equation}\label{ct:Lj}
 L_j=\frac{2^{2j}B_{2j+1}(1/4)}{j(2j+1)}\quad(j\ge1),
 \qquad L_1=\frac1{16},\quad L_2=-\frac5{128},\quad L_3=\frac{61}{768}.
\end{equation}
Define the formal series
\begin{equation}\label{ct:Tformal}
 T_s(v)=\exp\left(s\sum_{j\ge1}L_jv^j\right)
       \sum_{j\ge0}\frac{E_{2j}(s)_{2j}}{2^{2j}(2j)!}v^j
       =\sum_{j\ge0}t_j(s)v^j.
\end{equation}
Each coefficient is a finite rational polynomial; the infinite formal
series is not asserted to converge. An explicit terminating recursion is
\begin{align}
 e_0(s)&=1,&
 e_n(s)&=\frac{s}{n}\sum_{j=1}^n jL_je_{n-j}(s),\label{ct:recursion}\\
 t_n(s)&=\sum_{j=0}^ne_j(s)
       \frac{E_{2n-2j}(s)_{2n-2j}}{2^{2n-2j}(2n-2j)!}.\notag
\end{align}
In particular,
\begin{equation}\label{ct:first-t}
 t_0=1,\qquad t_1(s)=-\frac{s(2s+1)}{16},\qquad
 t_2(s)=\frac{s(20s^3+108s^2+211s+60)}{1536}.
\end{equation}

\begin{lemma}[Centered Gamma--Hurwitz expansion]\label{ct:centered}
For each $K\ge0$,
\begin{equation}\label{ct:tauexpansion}
 \tau_s(x+\tfrac12)=\frac12x^{-s/2}
       \sum_{j=0}^Kt_j(s)x^{-2j}
       +O\bigl(x^{-\Re s/2-2K-2}\bigr).
\end{equation}
The expansion is locally uniform in the order and in bounded complex
shifts of $x$ as $x$ tends to positive infinity. Fixed spectral derivatives
have the differentiated expansion, with finite powers of $\log x$;
fixed argument derivatives improve the decay by their order.
Furthermore $\deg_s t_j\le2j$.
\end{lemma}
\begin{proof}
The classical Gamma quotient expansion \cite{DLMF}, together with the alternating Hurwitz expansion \cite{HuKim}, gives
\[
 \log h(x+\tfrac12)=\tfrac12\log x+
                 \sum_{j=1}^K L_jx^{-2j}+O(x^{-2K-2}).
\]
Only even inverse powers remain because
$B_{n+1}(3/4)-B_{n+1}(1/4)=0$ for odd $n$; the surviving coefficient
is precisely \eqref{ct:Lj}. For $\Re s>0$ the alternating Hurwitz
integral is
\[
 \mathcal B_s(x+\tfrac12)=\frac1{2\Gamma(s)}
  \int_0^\infty t^{s-1}e^{-xt}\operatorname{sech}(t/2)\,dt.
\]
Expanding the smooth factor at zero gives
\[
 \mathcal B_s(x+\tfrac12)\sim\frac12x^{-s}
        \sum_{j\ge0}\frac{E_{2j}(s)_{2j}}{2^{2j}(2j)!}x^{-2j}.
\]
For an arbitrary compact order set the same finite expansion, with a
uniform remainder, follows by taking sufficiently many terms in the
ordinary Hurwitz Euler--Maclaurin expansions of \eqref{ct:B}. This avoids
using the displayed integral outside its convergence domain. Multiplying
finite expansions with their remainders proves \eqref{ct:tauexpansion}.
Cauchy estimates on slightly larger compact sets justify spectral
 derivatives; the differentiated Hurwitz and Gamma expansions justify
argument derivatives. Finally $\deg e_j\le j$, while
$\deg(s)_{2n-2j}=2n-2j$, proving the asserted degree bound.
\end{proof}

For an integer $k\ge0$ and a shift $v$ define
\begin{equation}\label{ct:bk}
 b_k(s,v)=\sum_{j=0}^{\lfloor k/2\rfloor}
 t_j(s)\binom{-s/2-2j}{k-2j}v^{k-2j}.
\end{equation}
Throughout, generalized binomial coefficients in such formulas mean their
polynomial expressions. They must not be computed by separately
specializing Gamma quotients at poles. We have
\begin{equation}\label{ct:b-properties}
 \deg_s b_k\le k,\qquad b_k(s,-v)=(-1)^kb_k(s,v).
\end{equation}

Choose a \emph{fixed} cutoff parameter $c$, initially independent of
$p,q,a,b$, and write
\begin{equation}\label{ct:fixed-cutoff}
 X=N+c,\quad \lambda=\frac{p+q}{2},\quad
 A=a+\tfrac12-c,\quad B=b+\tfrac12-c.
\end{equation}
The complete ambient coefficient polynomials are
\begin{align}
 c_k(p,q;a,b;c)
 &=\frac14\sum_{j=0}^k\bigl\{
 b_j(p,A)b_{k-j}(q,B+1)+b_j(p,A+1)b_{k-j}(q,B)\bigr\},
                         \label{ct:ck}\\
 d_k(p,q;a,b;c)
 &=\frac14\sum_{j=0}^k\bigl\{
 b_j(p,A)b_{k-j}(q,B+1)-b_j(p,A+1)b_{k-j}(q,B)\bigr\}.
                         \label{ct:dk}
\end{align}
Thus $c_0=1/2$, $d_0=0$, and both polynomials have total degree at most
$k$ in the independent orders $p,q$. Lemma~\ref{ct:centered} implies
\begin{align}
 C_N&=X^{-\lambda}\sum_{k=0}^Kc_kX^{-k}
          +O(X^{-\Re\lambda-K-1}),\label{ct:Cexp}\\
 D_N&=X^{-\lambda}\sum_{k=0}^Kd_kX^{-k}
          +O(X^{-\Re\lambda-K-1}).\label{ct:Dexp}
\end{align}
These expansions inherit the same uniformity and derivative properties.

\subsection{Every independent contact and every mixed order derivative}

Fix $p_0+q_0=-2m$, where $m\ge0$ is an integer. Differentiate at independent
ambient $p,q$ first, then specialize to $(p_0,q_0)$. The symbol
$\operatorname{FP}_X$ below denotes the coefficient of
$X^0(\log X)^0$ in the resulting asymptotic expansion.

\begin{theorem}[Complete independent contact calculus]\label{ct:contacts}
For every $r,s\ge0$,
\begin{align}
 \operatorname{FP}_X
 \left.\partial_p^r\partial_q^sS_N^+\right|_{(p_0,q_0)}
 &=\left.\partial_p^r\partial_q^sE_+\right|_{(p_0,q_0)}
    -\left.\partial_p^r\partial_q^sc_m\right|_{(p_0,q_0)}.
                                          \label{ct:pluscontact}
\end{align}
For the alternating law there are two parity finite parts. On the
subsequence $N\equiv\epsilon\pmod2$, using the same coordinate $X=N+c$,
\begin{align}
 \operatorname{FP}_{X,\epsilon}
 \left.\partial_p^r\partial_q^sS_N^-\right|_{(p_0,q_0)}
 &=\left.\partial_p^r\partial_q^sE_-\right|_{(p_0,q_0)}
   +(-1)^\epsilon\left.\partial_p^r\partial_q^sd_m\right|_{(p_0,q_0)}.
                                          \label{ct:minuscontact}
\end{align}
Every contact derivative of total order $r+s>m$ vanishes. The same formulas
hold after fixed argument derivatives if $c$ is held independent while
these derivatives are taken.
\end{theorem}
\begin{proof}
For any coefficient polynomial $a_k(p,q)$,
\begin{align}
 \partial_p^r\partial_q^s\{a_kX^{-\lambda-k}\}
 =X^{-\lambda-k}\sum_{i=0}^r\sum_{j=0}^s
 \binom ri\binom sj(\partial_p^i\partial_q^ja_k)
             \left(-\tfrac12\log X\right)^{r+s-i-j}.
 \label{ct:logjets}
\end{align}
After specialization, only $k=m$ has zero exponent in $X$, and only
$i=r,j=s$ has logarithmic degree zero. Choose $K$ large enough that the
remainder and all fixed derivatives tend to zero. Equations
\eqref{ct:finite}, \eqref{ct:Cexp}, and \eqref{ct:Dexp} prove both formulas.
The degree bound in \eqref{ct:ck}--\eqref{ct:dk} proves the last assertion
about spectral derivatives. Fixed-cutoff argument derivatives follow from
the same uniform expansions.
\end{proof}

The arithmetic mean of the two alternating parity finite parts is the
endpoint continuation, including all its mixed jets. A full-sequence
ordinary cutoff constant is not supplied when the two constants differ.
Re-expanding in the parity index $(N-\epsilon)/2$ would be a different
coordinate convention. Similarly, imposing $q=-2m-p$ before differentiating
would calculate only tangent derivatives; it would lose the independent
transverse derivatives required by Theorem~\ref{ct:contacts}.

\subsection{Balanced coordinates, low-order contacts, and exact parity}

The most symmetric cutoff is
\begin{equation}\label{ct:balanced}
 X=N+\mu,\qquad \mu=\frac{a+b}{2}+1,\qquad
 \delta=\frac{a-b}{2}.
\end{equation}
The four centered arguments in the two endpoint products are then
$X+\delta-1/2$, $X-\delta+1/2$, $X+\delta+1/2$, and
$X-\delta-1/2$. Abbreviate the corresponding polynomials to $c_k,d_k$,
and put $S=p+q$, $D=p-q$. The first ambient coefficients are
\begin{align}
 c_0&=\tfrac12,&d_0&=0,\notag\\
 c_1&=-\frac{\delta D}{4},&d_1&=\frac D8,\label{ct:lowambient}\\
 c_2&=\frac{4\delta^2(D^2+2S)-D^2-2S^2}{64},&
 d_2&=-\frac{\delta(D^2+2S)}{16}.\notag
\end{align}
The third coefficients, useful for independent jets at the next resonance,
are
\begin{align}
 c_3&=-\frac{\delta D}{384}
 \{4\delta^2(D^2+6S+8)-3D^2-6S^2-36S\},\label{ct:thirdambient}\\
 d_3&=\frac D{768}
 \{12\delta^2(D^2+6S+8)-5D^2-6S^2-48S-16\}.\notag
\end{align}
For zero-order contacts only, their restrictions to $S=-2m$ give the
following compact table. Here $d=a-b=2\delta$ and $D=p-q$.
\begin{center}
\renewcommand{\arraystretch}{1.6}
\begin{tabular}{c|c|c}
$m$&symmetric contact $c_m$&alternating amplitude $d_m$\\\hline
$0$&$1/2$&$0$\\
$1$&$-dD/8$&$D/8$\\
$2$&$\displaystyle\frac{(D^2-8)d^2-D^2-32}{64}$
   &$\displaystyle-\frac{d(D^2-8)}{32}$\\
$3$&$\displaystyle-\frac{dD\{(D^2-28)d^2-3D^2\}}{768}$
   &$\displaystyle\frac{D\{3d^2(D^2-28)-5D^2+56\}}{768}$
\end{tabular}
\end{center}
The table must not replace the ambient polynomials when calculating
independent spectral derivatives.

\begin{proposition}[Balanced parity]\label{ct:parity}
At the balanced cutoff,
\begin{align}
 c_k(p,q;-\delta)&=(-1)^kc_k(p,q;\delta),&
 d_k(p,q;-\delta)&=(-1)^{k+1}d_k(p,q;\delta),\label{ct:reflection}\\
 c_k(q,p;\delta)&=(-1)^kc_k(p,q;\delta),&
 d_k(q,p;\delta)&=(-1)^kd_k(p,q;\delta).\notag
\end{align}
Every odd symmetric contact is divisible by $(a-b)(p-q)$.
Every even alternating contact vanishes at equal shifts, and every odd
alternating contact vanishes at equal orders.
\end{proposition}
\begin{proof}
Use $b_k(s,-v)=(-1)^kb_k(s,v)$ in the two endpoint products. Reflecting
$\delta$ exchanges these products; interchanging the independent orders
then gives the second line. Polynomial divisibility follows from the
resulting sign changes.
\end{proof}

The first off-diagonal symmetric contact is therefore
$(q-p)(a-b)/8$ on $p+q=-2$, and it vanishes on either diagonal separately.
For example,
\[
 p=\tfrac13,\quad q=-\tfrac73,\quad a=\tfrac23,\quad b=\tfrac75
 \quad\Longrightarrow\quad c_1=\frac{11}{45},\quad d_1=\frac13.
\]
At $p=q=s$, $a=b$, one has the exact reduction
\[
 C_N=2(N+a+1)^s\mathcal B_s(N+a+1)\mathcal B_s(N+a+2),
 \qquad D_N=0.
\]
Consequently the present coefficients specialize to twice the source's
diagonal endpoint coefficients. This checks all its even-order contacts,
all vanishing odd contacts, and every diagonal spectral derivative.

\subsection{Translation, the sharp derivative threshold, and moving shifts}

If $Y=X-h$, coefficient transport is finite:
\begin{equation}\label{ct:translate}
 c_m^{[h]}(p,q)=\sum_{k=0}^m c_k(p,q)
      \binom{-\lambda-k}{m-k}h^{m-k},
 \qquad
 d_m^{[h]}(p,q)=\sum_{k=0}^m d_k(p,q)
      \binom{-\lambda-k}{m-k}h^{m-k}.
\end{equation}
These are ambient polynomial identities, obtained by expanding
$X^{-\lambda-k}=Y^{-\lambda-k}(1+h/Y)^{-\lambda-k}$.
They retain total spectral degree at most $m$. At zero derivative order
and on $\lambda=-m$, each binomial factor is one. For higher independent
jets it must first be differentiated as a polynomial. Translating a cutoff
therefore changes transverse contact derivatives even when the restricted
zero-order formula appears elementary.

There is also a closed expression for every highest mixed derivative.
For a polynomial $P(p,q)$ write $[P]_m$ for its homogeneous component of
total spectral degree $m$.

\begin{theorem}[Highest mixed contacts and sharpness]\label{ct:top}
For the balanced cutoff translated by $h$, the formal generating identities
are
\begin{align}
 \sum_{m\ge0}[c_m^{[h]}]_m t^m
 &=\frac12\exp\left(-\frac{(p-q)\delta+(p+q)h}{2}t\right)
 \frac{\cosh((p-q)t/4)}{\cosh(pt/2)\cosh(qt/2)},\label{ct:topplus}\\
 \sum_{m\ge0}[d_m^{[h]}]_m t^m
 &=\frac12\exp\left(-\frac{(p-q)\delta+(p+q)h}{2}t\right)
 \frac{\sinh((p-q)t/4)}{\cosh(pt/2)\cosh(qt/2)}.\label{ct:topminus}
\end{align}
The total spectral degree is generically exactly $m$ for $c_m$, and
exactly $m$ for $d_m$ when $m\ge1$. Hence the threshold $r+s>m$ in
Theorem~\ref{ct:contacts} is sharp.
\end{theorem}
\begin{proof}
In \eqref{ct:Tformal}, the highest spectral coefficient of $t_j(s)$ is
$E_{2j}s^{2j}/(2^{2j}(2j)!)$: every nonconstant contribution from the
Gamma exponential has strictly smaller degree. The highest-degree
series for a shifted $\tau_s$ is therefore
$\tfrac12 e^{-svt/2}\operatorname{sech}(st/2)$. Multiplying the
appropriate two factors in \eqref{ct:ck}--\eqref{ct:dk}, and taking
their sum or difference, proves \eqref{ct:topplus}--\eqref{ct:topminus}.
For $h=q=0$, the coefficient of $\delta^mp^m$ in the symmetric
expression is $(-1)^m/(2^{m+1}m!)$. For $m\ge1$, the coefficient
of $\delta^{m-1}p^m$ in the alternating expression is
$(-1)^{m-1}/(2^{m+2}(m-1)!)$. Both are nonzero.
\end{proof}

The balanced coordinate depends on $a,b$. To differentiate the original
finite sums in these shifts, use \eqref{ct:ck}--\eqref{ct:dk} with fixed
$c$, differentiate, and only then set $c=\mu$. Holding the balanced $X$
artificially fixed while differentiating the original $N$-cutoff would
omit its coordinate derivative. This restriction does not affect
independent order derivatives, because $\mu$ is independent of $p,q$.

\subsection{Ordinary polynomial and logarithmic subtractions}

The coefficient algebra gives convergent sums representing every contact
and every jet. It does not require assigning ordinary values to divergent
unmodified sums. For $\Re c>1$ and $K\ge0$ define
\begin{equation}\label{ct:PK}
 P_K^+(p,q;X)=\sum_{k=0}^Kc_kX^{-\lambda-k},\qquad
 P_K^-(p,q;X)=\sum_{k=0}^Kd_kX^{-\lambda-k},\qquad X_n=n+c.
\end{equation}

\begin{theorem}[Normally convergent continuation]\label{ct:normal}
For $\Re\lambda>-K-1$,
\begin{equation}\label{ct:normalplus}
 \sum_{n\ge0}\{g_n^+-P_K^+(X_n-1)+P_K^+(X_n)\}
             =E_+-P_K^+(c-1)
\end{equation}
converges absolutely and locally normally. For $\Re\lambda>-K$,
\begin{equation}\label{ct:normalminus}
 \sum_{n\ge0}(-1)^n\{g_n^--P_K^-(X_n-1)-P_K^-(X_n)\}
             =E_--P_K^-(c-1)
\end{equation}
also converges absolutely and locally normally. The latter is an ordinary
convergent series, normally convergent after pairing consecutive terms,
on the larger domain $\Re\lambda>-K-1$.
All fixed mixed parameter derivatives may be taken under the sum in the
stated absolute-normal domains. Formula~\eqref{ct:logjets} specifies
every polynomial and logarithmic counterterm explicitly.
\end{theorem}
\begin{proof}
Subtract the finite telescoping identities for $P_K^\pm$ from
\eqref{ct:finite}. The remaining endpoints are
$C_N-P_K^+(X_N)$ and $D_N-P_K^-(X_N)$, both
$O(X_N^{-\Re\lambda-K-1})$. They vanish on the larger asserted domain.
For the symmetric series, the summand is a unit difference of its endpoint
remainder, so its size is $O(n^{-\Re\lambda-K-2})$. For the alternating
series the sum of two adjacent remainders is
$O(n^{-\Re\lambda-K-1})$, giving its absolute domain. Pairing successive
alternating terms turns this sum into a difference over two steps and
restores the extra inverse power. The uniform differentiated expansions
and Cauchy estimates justify all fixed derivatives in the stated domains.
The condition $\Re c>1$ keeps the lower power $c-1$ on the fixed branch.
\end{proof}

At $\lambda=-m$, the symmetric subtraction with $K=m$ is a polynomial.
Writing
\begin{equation}\label{ct:resonantP}
 \Pi_m^+(X)=\sum_{k=0}^mc_kX^{m-k},\qquad
 \Pi_m^-(X)=\sum_{k=0}^md_kX^{m-k},
\end{equation}
we obtain the ordinary absolutely convergent identity
\begin{equation}\label{ct:polynomialplus}
 \sum_{n\ge0}\{g_n^+-\Pi_m^+(n+\mu-1)+\Pi_m^+(n+\mu)\}
              =E_+-\Pi_m^+(\mu-1).
\end{equation}
Its summand is $O(n^{-2})$. The polynomial alternating subtraction is
ordinary convergent, generally conditionally, with value
$E_--\Pi_m^-(\mu-1)$. Including the additional term
$d_{m+1}X^{-1}$ makes it absolutely convergent. Independent mixed jets
at the resonance use $K=m$ in the symmetric formula and $K=m+1$ in
the alternating formula, with the logarithmic counterterms
\eqref{ct:logjets}. Thus the complete contact theorem has ordinary
convergent-series representatives at every derivative order.

\subsection{The first off-diagonal resonance and exact Gamma examples}

\begin{corollary}[A short identity on $p+q=-2$]\label{ct:firstresonance}
For $p+q=-2$ and $\Re a,\Re b>0$,
\begin{equation}\label{ct:firstplus}
 \boxed{\displaystyle
 \sum_{n\ge0}\left\{
  [V_p(n+a+1)+V_q(n+b+1)]u_nv_n+\frac12\right\}
 =\tau_p(a)\tau_q(b+1)+\tau_p(a+1)\tau_q(b)+\frac{pa+qb}{4}.}
\end{equation}
The series is absolutely and locally normally convergent on this
hyperplane. There is also the absolutely convergent alternating identity
\begin{align}
 &\sum_{n\ge0}(-1)^n\left\{
 [V_p(n+a+1)-V_q(n+b+1)]u_nv_n-\frac{p-q}{4}\right.\notag\\
 &\hspace{18mm}\left.
 -\frac{pq(a-b)}8\left[\frac1{n+(a+b)/2}
                      +\frac1{n+(a+b)/2+1}\right]\right\}\notag\\
 &\qquad=\tau_p(a)\tau_q(b+1)-\tau_p(a+1)\tau_q(b)
            -\frac{p-q}{8}-\frac{pq(a-b)}{4(a+b)}.
                                                    \label{ct:firstminus}
\end{align}
\end{corollary}
\begin{proof}
Here $\Pi_1^+(X)=X/2+(q-p)(a-b)/8$ and
\[
 \Pi_1^+(\mu-1)=-\frac{pa+qb}{4},
\]
where the relation $p+q=-2$ is used only after coefficient extraction.
Equation~\eqref{ct:polynomialplus} proves \eqref{ct:firstplus}.
For the alternating formula, \eqref{ct:lowambient} gives
$d_1=(p-q)/8$ and $d_2=pq(a-b)/8$ at this resonance. Apply
\eqref{ct:normalminus} with $P_2^-(X)=d_1+d_2/X$.
\end{proof}

In particular take $p=1$, $q=-3$, $a=1$, $b=2$, and set
\[
 u_n=\tau_1(n+2),\qquad v_n=\tau_{-3}(n+3).
\]
The two exact evaluations are
\begin{align}
 \sum_{n\ge0}\{[V_1(n+2)+V_{-3}(n+3)]u_nv_n+\tfrac12\}
 &=\frac{55\pi\log2}{128}
      +\frac{9(1-\log2)}{4\pi}-\frac54,\label{ct:exampleplus}\\
 \sum_{n\ge0}(-1)^n\left\{
 [V_1(n+2)-V_{-3}(n+3)]u_nv_n-1\right.&\notag\\[-3pt]
 \left.\hspace{6mm}-\frac38\left(\frac1{n+3/2}+\frac1{n+5/2}\right)
 \right\}
 &=\frac{55\pi\log2}{128}
      -\frac{9(1-\log2)}{4\pi}-\frac34.\label{ct:exampleminus}
\end{align}
Both sums are ordinary and absolutely convergent. Their summands use only
Gamma quotients, digamma functions, rational functions, and an Euler
polynomial: $\mathcal B_1(z)=\tfrac12\{\psi((z+1)/2)-\psi(z/2)\}$ and
\[
 \mathcal B_{-3}(z)=\frac12E_3(z),\qquad
 E_3(z)=z^3-\frac32z^2+\frac14.
\]
For a direct evaluation of the endpoints, put $r=\sqrt{2/\pi}$. Then
$h(1)=r$, $h(2)=1/r$, $h(3)=2r$,
$\mathcal B_1(1)=\log2$, and $\mathcal B_1(2)=1-\log2$. Consequently
\[
 \tau_1(1)\tau_{-3}(3)=\frac{55\pi\log2}{128},\qquad
 \tau_1(2)\tau_{-3}(2)=\frac{9(1-\log2)}{4\pi},
\]
which proves both displayed constants.

\subsection{Which resonant unmodified sums actually converge?}

The contact polynomials also resolve this ordinary-convergence question,
which is stronger than asking whether one constant coefficient vanishes.

\begin{proposition}[Ordinary convergence at negative-total-order resonances]
\label{ct:resonantconvergence}
On $p+q=-2m$ with $m>0$, the unmodified symmetric series does not converge.
At $m=0$ it converges absolutely and has the inherited contact $1/2$.
The alternating series at $m=0$ always converges; its convergence is
conditional unless $p=q=0$, when every summand is zero.
For the unmodified alternating series at $m\ge1$ the following are exact:
\begin{enumerate}
 \item If $p\ne q$, it does not converge.
 \item If $p=q=-m$ and $a\ne b$, it converges when $m=1$ and does
       not converge when $m\ge2$.
 \item If $p=q$ and $a=b$, every alternating summand is zero.
\end{enumerate}
In item 2 at $m=1$ the convergence is generally conditional.
\end{proposition}
\begin{proof}
The symmetric endpoint has leading term $X^m/2$, which cannot tend to a
finite limit for $m>0$. At $m=0$ it tends to $1/2$, and its difference
is $O(n^{-2})$. At $m=0$, the alternating endpoint tends to zero;
its summand has leading amplitude $(p-q)/(4n)$ unless $p=q=0$.
In that exceptional case $V_0=0$ and the summands vanish.
If $m\ge1$ and $p\ne q$, the alternating endpoint has leading term
$(p-q)X^{m-1}/8$, so $(-1)^ND_N$ has no limit. If $p=q=-m$, parity
and \eqref{ct:lowambient} instead give
\[
 D_N=\frac{m(a-b)}8X^{m-2}+O(X^{m-4}).
\]
For distinct shifts this tends to zero exactly when $m=1$; for $m=2$
it yields two different parity limits, and for larger $m$ it grows.
The exact finite identity \eqref{ct:finite} proves convergence or its
failure in each case. At $m=1$ the alternating summand generally has
nonzero $1/n$ amplitude, establishing the final qualification. Equal
orders and shifts give $u_n=v_n$, hence $D_n=0$ identically.
\end{proof}

The theorem answers the source's contact question at all its specified
hyperplanes. Independent extra outer powers and longer products remain
separate summation problems. The present finite coefficients offer a
precise starting point for those extensions: they identify which endpoint
terms must be retained before any additional spectral differentiation.
